% Contenido del póster: Fareed y Viotto (2025), Finance Research Letters. DOI 10.1016/j.frl.2024.106642
% Elaborado con apoyo de IA a partir del resumen oficial del artículo (OpenAlex), 2026-11-12.
\def\DOI{10.1016/j.frl.2024.106642}
\def\TituloEs{Plataformas de crowdfunding e inclusión financiera: ¿promesa cumplida o desilusión?}
\def\TituloOriginal{Crowdfunding platforms and financial inclusion: Fulfilled promise or disillusion?}
\def\Autores{Fareed y Viotto}
\def\Anio{2025}
\def\Revista{Finance Research Letters}

\def\Problema{Las plataformas de crowdfunding con fines sociales prometen ampliar la inclusión financiera en los países en desarrollo, pero la literatura no tiene evidencia de si cumplen esa promesa ni en qué medida.}
\def\Objetivo{Evaluar si las plataformas de crowdfunding con fines sociales mejoran la inclusión financiera en los países en desarrollo.}
\def\Metodo{Los autores aplican distintas técnicas para tratar la endogeneidad entre la actividad de las plataformas y los resultados de inclusión.}
\def\Resultados{\begin{itemize}
  \item Las plataformas amplían el acceso al capital a más prestatarios.
  \item La atención a las personas más vulnerables no mejora.
  \item Las tasas de interés no cambian.
\end{itemize}}
\def\Conclusion{La promesa se cumple solo en parte: las plataformas llegan a más prestatarios, pero no a los más vulnerables ni a menores costos. El estudio ofrece recomendaciones para las políticas que vinculan tecnologías digitales e inclusión financiera.}
\def\Relevancia{En el Perú, el financiamiento participativo también amplió su alcance, pero hacia proyectos empresariales cada vez más grandes: el monto medio por proyecto pasó de S/~\cifra{montoMedio2022} mil en 2022 a S/~\cifra{montoMedio2026} mil en 2026, y el segmento personal desapareció tras \cifra{finTercera}. Como en este estudio, crecer no equivale a llegar a los prestatarios más pequeños.}
\def\PalabrasClave{crowdfunding, inclusión financiera, países en desarrollo, tasas de interés, plataformas sociales}
% --- Ficha del estudio (herramientas/posteres/ficha_articulos.py) ---
\def\Ambito{Países en desarrollo}
\def\Enfoque{Empírico, con técnicas para tratar la endogeneidad}
\def\FiguraPeru{2}
% --- fin de la ficha ---
